\documentclass[a4paper,12pt]{article}
\usepackage[T2A]{fontenc}
\usepackage[utf8]{inputenc}
\usepackage[russian]{babel}
\usepackage{amsmath,amssymb,mathtools}
\usepackage{geometry}
\usepackage{booktabs,tabularx,array}
\usepackage{enumitem}
\usepackage{microtype}
\usepackage{tikz}
\usepackage{xcolor}
\usepackage{hyperref}
\usepackage{fancyhdr}
\usepackage{parskip}
\geometry{left=22mm,right=22mm,top=23mm,bottom=25mm}
\renewcommand{\familydefault}{\rmdefault}
\setlength{\headheight}{15pt}
\definecolor{accent}{HTML}{007E80}
\definecolor{ink}{HTML}{262F35}
\definecolor{errorred}{HTML}{A42B30}
\definecolor{keyblue}{HTML}{225D7A}
\definecolor{softgray}{HTML}{58656B}
\hypersetup{colorlinks=true,linkcolor=accent,urlcolor=accent}
\setlength{\parindent}{0pt}
\setlength{\parskip}{6pt}
\setlist[itemize]{leftmargin=*,nosep,topsep=3pt}
\pagestyle{fancy}
\fancyhf{}
\fancyhead[L]{\small\color{softgray}Математика \enspace/\enspace Ревью домашней работы}
\fancyhead[R]{\small\color{softgray}05.11.2026}
\fancyfoot[C]{\small\color{softgray}\thepage}
\renewcommand{\headrulewidth}{0pt}
\renewcommand{\footrulewidth}{0pt}
\newcommand{\labeltext}[1]{\par\noindent\textbf{\color{ink}#1}\quad}
\newcommand{\err}{\textbf{\color{errorred}Ошибка:}\ }
\newcommand{\idea}{\textbf{\color{keyblue}Ключевая идея:}\ }
\newcommand{\result}{\textbf{\color{accent}Итог:}\ }
\begin{document}
\thispagestyle{empty}
\begin{center}
\vspace*{34mm}
{\color{accent}\large МАТЕМАТИКА}\par
\vspace{19mm}
{\fontsize{24}{29}\selectfont\color{ink}Ревью домашней работы}\par
\vspace{12mm}
{\large\color{accent}Модуль и область определения функции}\par
\vspace{5mm}
{\normalsize Разбор заданий 2 и 3}\par
\vfill
{\normalsize 5 ноября 2026 г.}\par
\vspace{12mm}
{\normalsize Лёвин Артём Александрович}\par
{\small эксклюзивно для Анны Трапезниковой}\par
\vspace{27mm}
\begin{tikzpicture}[x=1cm,y=1cm]
\draw[accent,line width=.95pt] (0,0) .. controls (1.7,.7) and (3.3,-.48) .. (4.7,0)
 .. controls (6.2,.58) and (7.7,-.44) .. (9,0);
\end{tikzpicture}
\vspace{9mm}
\end{center}
\clearpage

\begin{center}
{\LARGE\color{ink}Сводка ошибок}\par
\vspace{7mm}
\end{center}
\small
\renewcommand{\arraystretch}{1.38}
\begin{tabularx}{\linewidth}{@{}p{11mm}p{27mm}X p{21mm}p{27mm}@{}}
\toprule
\textbf{№} & \textbf{Тема} & \textbf{Суть ошибки} & \textbf{Ваш ответ} & \textbf{Верный ответ}\\
\midrule
\hyperref[task:2]{2} & Минимум суммы модулей & Значение найдено в одной точке; нижняя оценка и все точки минимума не обоснованы. & $2$ & $2$ при $x\in[1;3]$\\
\addlinespace
\hyperref[task:3]{3} & Область определения корня & Модуль раскрыт без учёта знака; неверно определены промежутки и сумма целых чисел. & $11$ & $5$\\
\bottomrule
\end{tabularx}
\normalsize
\vspace{8mm}
\textcolor{keyblue}{\textbf{Основное внимание:}} для задачи на минимум требуется доказать, что меньших значений нет; в задачах с модулем нужно заранее найти все точки смены знака.
\clearpage

\thispagestyle{empty}
\vspace*{.34\textheight}
\begin{center}
{\Huge\color{accent}Разбор ошибок}
\end{center}
\clearpage

\phantomsection\label{task:2}
\section*{Задание 2. Минимум суммы модулей}
\labeltext{Условие.}
Докажите, что минимум функции
\[
f(x)=|x-3|+|x-1|,\qquad x\in[-1;3],
\]
составляет $2$.

\labeltext{Ваше решение.}
При $x=1$ вычислено:
\[
f(1)=|1-3|+|1-1|=2+0=2.
\]
После этого сделан вывод о минимальном значении функции.

\labeltext{Ваш ответ.} $2$.

\labeltext{Ошибка.} \err Недостаточно доказано, что найденное значение является наименьшим.

\labeltext{Где именно возникает ошибка.}
Последний правильный шаг --- вычисление $f(1)=2$. Первый логически неверный переход --- вывод «минимум равен $2$» непосредственно из подстановки $x=1$.

\labeltext{Почему этот переход недостаточен.}
Значение функции в одной точке показывает лишь, что число $2$ \emph{достигается}. Для доказательства минимума требуется ещё показать, что при каждом допустимом $x$ выполнено $f(x)\ge 2$. Например, само по себе равенство $h(1)=2$ ничего не говорит о том, может ли $h(0)$ оказаться меньше $2$.

\labeltext{Нарушенное правило.}
Для утверждения «$m$ --- минимум функции на заданном множестве» нужны два факта:
\begin{enumerate}[leftmargin=*,nosep]
\item $f(x)\ge m$ для всех допустимых $x$;
\item существует допустимая точка $x_0$, для которой $f(x_0)=m$.
\end{enumerate}

\labeltext{Ключевая идея.}
\idea Числа $|x-1|$ и $|x-3|$ --- расстояния от точки $x$ на числовой прямой до точек $1$ и $3$. Если $x$ лежит между $1$ и $3$, сумма этих расстояний равна расстоянию между крайними точками: $3-1=2$. Вне отрезка она больше.

\labeltext{Исправление Вашего решения.}
Сохраняем найденное равенство $f(1)=2$. Добавляем обоснование для \emph{всех} $x\in[-1;3]$. Проверим два промежутка, на которых знаки выражений под модулями постоянны.

\labeltext{Полное правильное решение.}
\textbf{1.} Если $-1\le x\le1$, то $x-3\le0$ и $x-1\le0$. Поэтому
\[
\begin{aligned}
f(x)&=(3-x)+(1-x)=4-2x.
\end{aligned}
\]
Так как $x\le1$, имеем $-2x\ge-2$, а значит,
\[
f(x)=4-2x\ge 4-2=2.
\]
Равенство на этом промежутке достигается при $x=1$.

\textbf{2.} Если $1\le x\le3$, то $x-3\le0$, а $x-1\ge0$. Следовательно,
\[
\begin{aligned}
f(x)&=(3-x)+(x-1)\\
&=3-x+x-1=2.
\end{aligned}
\]
На \emph{всём} отрезке $[1;3]$ функция равна $2$.

\textbf{3.} Отрезки $[-1;1]$ и $[1;3]$ вместе образуют всю заданную область $[-1;3]$. Значит, для любого допустимого $x$ выполняется $f(x)\ge2$, причём значение $2$ достигается.

\labeltext{Проверка.}
\[
f(-1)=|-4|+|-2|=6,\quad f(1)=2,\quad f(2)=1+1=2,\quad f(3)=0+2=2.
\]
Эти значения согласуются с доказанным результатом.

\labeltext{Итог.} \result
\[
\boxed{\min\limits_{x\in[-1;3]} f(x)=2},
\qquad \boxed{x\in[1;3]\text{ --- точки минимума}.}
\]

\labeltext{Задача на закрепление.}
Докажите, что минимум функции
\[
g(x)=|x+2|+|x-4|,\qquad x\in[-5;6],
\]
равен $6$. Найдите все значения $x$, при которых он достигается.
\clearpage

\phantomsection\label{task:3}
\section*{Задание 3. Область определения радикала}
\labeltext{Условие.}
Найдите сумму всех целых чисел, принадлежащих области определения функции
\[
f(x)=\sqrt{\bigl(|x-5|-3\bigr)\bigl(24-x^2+5x\bigr)}.
\]

\labeltext{Ваше решение.}
Верно записано требование к подкоренному выражению:
\[
\bigl(|x-5|-3\bigr)\bigl(24-x^2+5x\bigr)\ge0.
\]
Затем вместо $|x-5|-3$ записано $x-5-3=x-8$ без разбора случаев. Найдены значения $x=8$ и $x=-3$, построена схема знаков с промежутком от $-3$ до $8$; указан ответ $11$.

\labeltext{Ваш ответ.} $11$.

\labeltext{Ошибка.} \err Модуль раскрыт как $|x-5|=x-5$ для всех $x$; это верно только при $x\ge5$.

\labeltext{Где именно возникает ошибка.}
Последний правильный шаг:
\[
\bigl(|x-5|-3\bigr)\bigl(24-x^2+5x\bigr)\ge0.
\]
Первый неверный переход:
\[
|x-5|-3\ \longrightarrow\ x-5-3=x-8
\quad\text{без условия }x\ge5.
\]
Например, при $x=0$ исходный множитель равен $|{-5}|-3=2$, тогда как подставленный множитель равен $-8$. Даже знаки получились разными.

\labeltext{Почему это неверно.}
Модуль раскрывают по знаку выражения внутри него:
\[
|x-5|=
\begin{cases}
5-x,&x<5,\\
x-5,&x\ge5.
\end{cases}
\]
Поэтому
\[
|x-5|-3=
\begin{cases}
2-x,&x<5,\\
x-8,&x\ge5.
\end{cases}
\]
Именно поэтому у первого множителя \emph{два} нуля, $x=2$ и $x=8$. При исследовании только уравнения $x-8=0$ точка $x=2$ теряется.

\labeltext{Нарушенное правило.}
Знак произведения определяется знаками \emph{всех} множителей. Сначала нужно найти все их нули, а затем исследовать знаки на промежутках, разделённых этими точками. Нули подкоренного выражения допустимы: $\sqrt0=0$.

\labeltext{Ключевая идея.} \idea
Для квадратного корня требуется неотрицательное подкоренное выражение. Разложим второй множитель на линейные и исследуем знаки двух множителей с границами $-3$, $2$ и $8$.

\labeltext{Исправление от последнего правильного шага.}
\[
\begin{aligned}
A(x)&=|x-5|-3,\\
B(x)&=24-x^2+5x=(8-x)(x+3).
\end{aligned}
\]
Решая $A(x)=0$, получаем $|x-5|=3$, то есть
\[
x-5=-3\ \text{или}\ x-5=3
\quad\Longrightarrow\quad x=2\ \text{или}\ x=8.
\]
Из $B(x)=0$ получаем $x=-3$ или $x=8$.

\labeltext{Полное правильное решение.}
Для $A(x)=|x-5|-3$:
\begin{itemize}
\item $A(x)>0$ при $x<2$ или $x>8$;
\item $A(x)<0$ при $2<x<8$;
\item $A(x)=0$ при $x=2$ и $x=8$.
\end{itemize}
Для $B(x)=(8-x)(x+3)$:
\begin{itemize}
\item $B(x)>0$ при $-3<x<8$;
\item $B(x)<0$ при $x<-3$ или $x>8$;
\item $B(x)=0$ при $x=-3$ и $x=8$.
\end{itemize}
Сопоставим знаки на промежутках:
\begin{center}
\small
\renewcommand{\arraystretch}{1.22}
\begin{tabular}{@{}lcccc@{}}
\toprule
$x$ & $(-\infty;-3)$ & $(-3;2)$ & $(2;8)$ & $(8;+\infty)$\\
\midrule
$A(x)$ & $+$ & $+$ & $-$ & $+$\\
$B(x)$ & $-$ & $+$ & $+$ & $-$\\
$A(x)B(x)$ & $-$ & $+$ & $-$ & $-$\\
\bottomrule
\end{tabular}
\end{center}
В точках $x=-3$, $x=2$, $x=8$ хотя бы один множитель обращается в нуль. Поэтому все три точки допустимы.

Следовательно, область определения функции:
\[
\boxed{D(f)=[-3;2]\cup\{8\}}.
\]
Целые числа в ней:
\[
-3,\ -2,\ -1,\ 0,\ 1,\ 2,\ 8.
\]
Их сумма равна
\[
(-3)+(-2)+(-1)+0+1+2+8=5.
\]

\labeltext{Проверка.}
При $x=0$ подкоренное выражение равно
\[
(|-5|-3)\cdot24=2\cdot24=48>0.
\]
При $x=4$ оно равно
\[
(|-1|-3)\cdot(24-16+20)=(-2)\cdot28=-56<0.
\]
При $x=8$ оба множителя равны нулю, значит, точка допустима. Также сумма всех целых чисел от $-3$ до $8$ \emph{без пропусков} составила бы $30$, а не $11$; это дополнительный способ заметить ошибку в подсчёте.

\labeltext{Итог.} \result
\[
\boxed{-3-2-1+0+1+2+8=5}.
\]

\labeltext{Задача на закрепление.}
Найдите сумму всех целых чисел, принадлежащих области определения функции
\[
g(x)=\sqrt{\bigl(|x-3|-2\bigr)\bigl(12+x-x^2\bigr)}.
\]

\vfill
\begin{center}
{\footnotesize\color{softgray}Лёвин Артём Александрович эксклюзивно для Анны Трапезниковой}
\end{center}
\end{document}
